\documentclass[11pt,a4paper]{article}

\usepackage[utf8]{inputenc}
\usepackage[T1]{fontenc}
\usepackage[provide=*,spanish]{babel}
\usepackage{lmodern}
\usepackage{microtype}
\usepackage[a4paper,margin=2.5cm]{geometry}
\usepackage{booktabs}
\usepackage{tabularx}
\usepackage{longtable}
\usepackage{array}
\usepackage{enumitem}
\usepackage{xcolor}
\usepackage[round,authoryear]{natbib}
\usepackage{hyperref}

\definecolor{azulieb}{HTML}{174A6E}
\definecolor{grisclaro}{HTML}{F1F4F6}
\hypersetup{
  colorlinks=true,
  linkcolor=azulieb,
  urlcolor=azulieb,
  pdftitle={Marco Curricular para la Formación en Informática de la Biodiversidad en la Era de los Agentes de IA},
  pdfsubject={Outline de manuscrito para el taller internacional, septiembre de 2026}
}

\setcounter{secnumdepth}{2}
\setcounter{tocdepth}{2}
\setlist[itemize]{leftmargin=1.5em,itemsep=0.25em,topsep=0.4em}
\setlist[enumerate]{leftmargin=1.7em,itemsep=0.25em,topsep=0.4em}
\renewcommand{\arraystretch}{1.2}

\newcommand{\pendiente}[1]{\textcolor{azulieb}{\textbf{[Por acordar:} #1\textbf{]}}}
\newenvironment{preguntas}
  {\par\medskip\begin{quote}\color{azulieb}\hrule\smallskip
   \textbf{Preguntas para el consenso}\par\smallskip\begin{itemize}}
  {\end{itemize}\smallskip\hrule\end{quote}\par\medskip}

\title{\textbf{Marco Curricular para la Formación en Informática de la
Biodiversidad en la Era de los Agentes de Inteligencia Artificial}\\[0.6em]
\large Outline de manuscrito para escritura colaborativa}
\author{Taller Internacional\\
Instituto de Ecología y Biodiversidad (Chile)\\
Instituto de Investigación de Recursos Biológicos Alexander von Humboldt (Colombia)\\[0.5em]
\pendiente{autorías y afiliaciones}}
\date{Concepción, Chile --- 14 al 17 de septiembre de 2026}

\begin{document}

\maketitle

\begin{abstract}
\noindent
La irrupción de agentes de inteligencia artificial capaces de generar código, consultar literatura, 
integrar herramientas y coordinar flujos de análisis está transformando la informática de la biodiversidad y 
cuestionando los modelos formativos centrados en lenguajes de programación y tutoriales específicos. Paralelamente, 
la implementación del Marco Global de Biodiversidad Kunming–Montreal exige capacidades sostenidas para integrar 
observaciones heterogéneas, producir indicadores y convertirlos en evidencia trazable para la toma de decisiones. 
Este artículo propone un marco curricular latinoamericano destinado a preparar estudiantes, investigadores, profesionales 
y equipos institucionales para colaborar responsablemente con agentes de IA en investigación, monitoreo y gestión de la 
biodiversidad. A partir de una síntesis conceptual de la literatura, el marco distingue diferentes grados de autonomía y 
vincula su uso con competencias en formulación de problemas, fundamentos computacionales, datos de biodiversidad, 
reutilización de repositorios, control de versiones, supervisión de agentes, reproducibilidad, validación, ética, gobernanza 
y colaboración interdisciplinaria. El modelo pedagógico combina problemas auténticos, práctica manual, asistencia de agentes, 
auditoría colectiva y evaluación de procesos, decisiones e incertidumbres, además de los productos finales. Se propone una 
arquitectura modular y adaptable, acompañada por un repositorio abierto y una comunidad regional de práctica que sostengan 
su actualización, reutilización y gobernanza. La tesis central es que la automatización no elimina la necesidad de fundamentos 
técnicos y ecológicos, sino que desplaza el énfasis desde la producción manual de código hacia la capacidad de comprender, 
reutilizar, auditar y gobernar ecosistemas abiertos de conocimiento, preservando la responsabilidad humana, la capacidad local y 
la cooperación científica.


\medskip
\noindent\textbf{Palabras clave:} informática de la biodiversidad; agentes de
inteligencia artificial; formación permanente; ciencia abierta; reproducibilidad;
auditoría; control de versiones; infraestructura digital.
\end{abstract}

\tableofcontents
\newpage

\section*{Propósito y modo de uso del outline}
\addcontentsline{toc}{section}{Propósito y modo de uso del outline}

Este documento es la estructura de trabajo del producto principal del taller:
un marco curricular que pueda orientar cursos, programas de capacitación,
escuelas de verano y actividades de fortalecimiento de capacidades en América
Latina. Cada sección corresponde a un bloque del programa y debe cerrarse con
acuerdos, disensos relevantes, evidencia o ejemplos y recomendaciones
accionables.

\begin{table}[ht]
\centering
\small
\begin{tabularx}{\textwidth}{@{}p{2.2cm}p{2.7cm}X p{3.0cm}@{}}
\toprule
\textbf{Bloque} & \textbf{Eje} & \textbf{Pregunta central} &
\textbf{Producto escrito}\\
\midrule
Lunes AM & Contexto &
¿Qué capacidades informáticas requiere la implementación del marco de monitoreo
del KMGBF y por qué la IA obliga a replantear la formación? & Introducción\\
Lunes PM & Transformación &
¿Cómo cambian los agentes la investigación, el desarrollo, la innovación, el
monitoreo y la gestión de la biodiversidad? & Sección 1\\
Martes AM & Competencias &
¿Qué conocimientos, habilidades y disposiciones debemos enseñar? & Sección 2\\
Martes PM & Pedagogía &
¿Cómo deberían enseñarse y evaluarse estas competencias? & Sección 3\\
Miércoles AM & Currículo &
¿Cómo estructurar un programa moderno, progresivo y adaptable? & Sección 4\\
Miércoles PM & Implementación &
¿Qué infraestructura, gobernanza y comunidad se requieren? & Sección 5\\
Jueves AM & Integración &
¿Qué mensajes, recomendaciones y compromisos emergen? & Conclusiones y hoja de ruta\\
\bottomrule
\end{tabularx}
\caption{Correspondencia entre los bloques del taller y el manuscrito.}
\end{table}

\clearpage
\section*{Charlas de provocación}
\addcontentsline{toc}{section}{Charlas de provocación}

Las charlas de provocación aportan un lenguaje común y preguntas iniciales para
la discusión y la escritura colaborativa. Sus ideas deberán registrarse como
insumos del taller y contrastarse con la literatura y la experiencia de las
personas participantes antes de incorporarlas como acuerdos del manuscrito. Las
intervenciones sin horario asignado se mantienen como ``por programar'' hasta
actualizar el programa cronológico.

\begin{center}
\footnotesize
\begin{tabularx}{\textwidth}{@{}
  >{\raggedright\arraybackslash}p{2.2cm}
  >{\raggedright\arraybackslash}p{3.4cm}
  >{\raggedright\arraybackslash}X
  >{\raggedright\arraybackslash}p{3.5cm}@{}}
\toprule
\textbf{Día y horario} & \textbf{Expositor(es)} &
\textbf{Provocación o eje} & \textbf{Secciones que alimenta}\\
\midrule
Lunes 14, 09:00--10:30 &
Diana Pulido, Instituto Humboldt (remoto) &
Contexto de biodiversidad y fortalecimiento de capacidades. &
Introducción; contexto del KMGBF; capacidades para investigación, monitoreo y
gestión.\\
Lunes 14, 10:50--12:00 &
Carlos Gortaris (remoto); Diego Arenas, DFKI (remoto) &
Agentes de IA: oportunidades, riesgos y transformaciones para la biodiversidad. &
Transformación de la I+D+i; riesgos y límites; colaboración humano--agente.\\
Por programar &
Leonardo Buitrago, Consultor Informática de la Biodiversidad &
Datos de biodiversidad: fuentes, calidad, integración, estándares y
procedencia. &
Marco de competencias; datos de biodiversidad; infraestructura y
reproducibilidad.\\
Por programar &
José Díaz, Boquila &
¿Qué se debe saber para trabajar en biodiversidad usando agentes de IA? &
Marco de competencias; resultados de aprendizaje; progresión curricular.\\
Por programar &
Luisa Diele-Viegas, ABECO &
Dimensiones sociales de la ciencia que podrían verse afectadas por la irrupción
de los agentes de inteligencia artificial. &
Reorganización del trabajo científico; riesgos; colaboración, ética y
gobernanza.\\
Martes 15, 14:00--15:30 &
Mario Miqueles, Creatividad y Autoformación &
Creatividad y autoformación: autoaprendizaje y modelos pedagógicos. &
Modelo pedagógico; actividades de aprendizaje; evaluación y aprendizaje
permanente.\\
Por programar &
Valeria Bermedo, Formación Permanente UdeC &
Elementos de diseño curricular para programas de formación permanente. &
Modelo pedagógico; arquitectura curricular; trayectorias, formatos y
evaluación.\\
\bottomrule
\end{tabularx}
\par\medskip
\textit{Charlas de provocación, expositores y contribución prevista al outline.}
\end{center}

\section{Introducción: por qué es necesario un nuevo marco curricular}

\subsection{Contexto de biodiversidad y fortalecimiento de capacidades}

\begin{itemize}
  \item Situar la propuesta ante la pérdida acelerada de biodiversidad y el
  incumplimiento global de las metas de Aichi para 2020, antecedentes que
  impulsaron la adopción del Marco Global de Biodiversidad Kunming--Montreal
  (KMGBF) del Convenio sobre la Diversidad Biológica. Sus cuatro objetivos para
  2050 y 23 metas de acción para 2030 vinculan la ambición global con estrategias,
  monitoreo y reportes nacionales; su implementación exige capacidades sostenidas
  para investigación, observación, gestión de datos, evaluación y toma de
  decisiones \citep{affinito2024monitoring,affinito2025coverage}.
  \item Describir brevemente las brechas regionales de formación, infraestructura
  y continuidad institucional que motivan una propuesta latinoamericana de
  formación permanente. Un análisis de 74,6 millones de publicaciones y 4,2
  millones de programas de curso encontró que, fuera de las disciplinas
  computacionales, la oferta formativa en IA no se corresponde bien con su uso y
  beneficio potencial en investigación. Aunque sus datos terminan en 2019 y sus
  medidas son aproximaciones bibliométricas, el estudio ofrece evidencia de una
  brecha entre adopción y formación que este currículo busca atender
  \citep{gao2024quantifying}.
 \item Caracterizar la investigación en biodiversidad como un dominio espacial, 
 temporal y multiescalar que integra ocurrencias y abundancias de especies, monitoreo 
 imperfecto, teledetección, presiones antropogénicas, clima, costos y conocimiento ecológico. 
 Dentro de este marco deben distinguirse dos ámbitos conectados: la investigación, que formula 
 y contrasta preguntas para producir conocimiento, y el monitoreo, una actividad aplicada, 
 sostenida y orientada a decisiones, que establece líneas base y repite observaciones comparables 
 para estimar el estado de la biodiversidad, sus cambios y la efectividad de las acciones. El monitoreo 
 requiere definir por qué, qué y cómo observar, además de sostener capacidades taxonómicas, series
 temporales, infraestructura y vínculos entre quienes producen y utilizan los datos \citep{paton2009biodiversity,buckland2005monitoring}. 
 Dos aplicaciones recientes de aprendizaje por refuerzo muestran el potencial de integrar estas dimensiones.
  \item Precisar los públicos del marco: estudiantes, investigadores,
  profesionales, gestores, desarrolladores, docentes y equipos institucionales.
\end{itemize}

\subsection{El cambio tecnológico}

\begin{itemize}
  \item Contrastar el modelo tradicional, centrado en tutoriales, lenguajes y
  análisis específicos, con un escenario en el que agentes generan código,
  construyen flujos de trabajo, interpretan literatura y asisten múltiples etapas
  de la investigación. La literatura describe este cambio como una transición
  desde herramientas expertas para tareas discretas hacia asistentes con autonomía
  parcial y, todavía de manera emergente, socios capaces de intervenir en ciclos
  completos de descubrimiento \citep{wei2025agentic}.
  \item Explicar el desplazamiento de valor desde la escritura manual de código
  hacia la formulación de problemas, comprensión de sistemas, reutilización,
  validación, mantenimiento y colaboración. El objetivo no debería ser la
  automatización completa de la ciencia, sino una colaboración
  \emph{con el piloto al mando}: equipos de agentes especializados pueden
  acelerar y coordinar tareas, mientras las personas conservan la autoridad y la
  responsabilidad de formular las preguntas, validar el proceso y aprobar las
  conclusiones. Este modelo exige formar científicos capaces de dirigir,
  cuestionar y auditar a los agentes, y de sostener la trazabilidad, la
  reproducibilidad y la rendición de cuentas del trabajo científico
  \citep{wang2026five}.
  \item Mantener como capacidad basal una experiencia mínima de programación
  manual, indispensable para leer, probar y auditar lo producido con agentes. La
  automatización no elimina la fragilidad técnica: los agentes aún pueden producir
  planes no ejecutables, código incorrecto y trayectorias muy sensibles a la
  formulación de las instrucciones \citep{wei2025agentic}.
\end{itemize}

\subsection{Problema, alcance y objetivo del manuscrito}

\begin{itemize}
  \item Formular el problema curricular: ¿cómo preparar personas y comunidades
  para colaborar con agentes sin perder comprensión crítica, autonomía,
  reproducibilidad ni cooperación humana?
  \item Delimitar el alcance a informática aplicada a investigación, monitoreo y
  gestión de la biodiversidad, incluyendo datos, software, modelos,
  infraestructura y prácticas comunitarias.
  \item Enunciar el objetivo general: proponer un currículo regional que permita
  reutilizar, comprender, evaluar, mantener y compartir ecosistemas abiertos de
  conocimiento.
  \item Presentar las tres preguntas rectoras: qué enseñar, cómo enseñarlo y cómo
  evaluarlo cuando humanos y agentes producen conocimiento conjuntamente.
\end{itemize}

\begin{preguntas}
  \item ¿Qué necesidades del marco de monitoreo del KMGBF deben quedar
  explícitamente vinculadas al currículo?
  \item ¿Cuáles son las brechas regionales prioritarias y qué evidencia o casos
  permiten describirlas?
  \item ¿Qué queda dentro y fuera del alcance de esta primera versión?
\end{preguntas}

\noindent\textbf{Resultado esperado de la sección:}
una tesis compartida sobre la necesidad, propósito, alcance y audiencia del marco
curricular. \pendiente{2--3 mensajes centrales y referencias de respaldo}

\section{Transformación de la I+D+i y de la gestión de la biodiversidad}

\subsection{Una escala para distinguir grados de autonomía}

No toda aplicación de IA constituye un agente ni transfiere la misma
responsabilidad. Para ordenar la discusión, se propone adaptar la escala de
\citet{wei2025agentic}:

\begin{enumerate}
  \item \textbf{Herramienta experta:} resuelve una tarea acotada; la formulación,
  ejecución e interpretación permanecen bajo control humano.
  \item \textbf{Asistente de investigación:} ejecuta autónomamente una etapa o
  subobjetivo predefinido y devuelve el control a la persona.
  \item \textbf{Socio científico autónomo:} recorre varias etapas del ciclo,
  formula hipótesis y adapta su estrategia con intervención humana reducida.
  \item \textbf{Arquitecto generativo:} inventaría nuevos instrumentos, métodos o
  marcos conceptuales; es una posibilidad prospectiva, no una capacidad que deba
  presumirse en sistemas actuales.
\end{enumerate}

Esta escala debe usarse para decidir qué tareas delegar, qué evidencia exigir y
qué supervisión corresponde a cada nivel, sin equiparar fluidez lingüística con
autonomía científica efectiva.

\noindent\textbf{Distinción operacional:} la IA generativa produce artefactos a
partir de instrucciones, mientras la IA agentiva combina razonamiento iterativo
---incluidas planificación y reflexión--- con interacción sostenida con
herramientas o entornos para perseguir objetivos
\citep{schneider2025generative}. Por tanto, una respuesta elaborada de un modelo
no basta para clasificarlo como agente.

\subsection{Capacidades emergentes y oportunidades}

\citet{wei2025agentic} sintetizan cinco capacidades que sostienen la agencia
científica: planificación y razonamiento, integración de herramientas, memoria,
colaboración entre agentes y optimización mediante retroalimentación. Estas
capacidades pueden combinarse a lo largo de cuatro procesos: observación y
formulación de hipótesis; planificación y ejecución; análisis de datos y
resultados; y síntesis, validación y evolución. De este modo en el campo de la
informática de la biodiversidad, los agentes pueden ofrecer oportunidades para:

\begin{itemize}
  \item Generación y revisión asistida de código y documentación.
  \item Exploración y síntesis de literatura científica.
  \item Construcción y orquestación de flujos de datos y análisis.
  \item Acceso en lenguaje natural a repositorios, servicios, APIs y servidores
  del Model Context Protocol (MCP).
  \item Apoyo al monitoreo aplicado mediante agentes que ayuden a coordinar
  adquisición desde sensores, colecciones y ciencia participativa; armonizar
  taxonomías y metadatos; detectar anomalías; estimar indicadores y su
  incertidumbre; mantener series temporales; y generar alertas e informes
  trazables  \citep{paton2009biodiversity,buckland2005monitoring,feng2022databases}.
  \item Expansión hacia preguntas y campos nuevos. Antes del auge reciente de los
  LLM, la utilización de IA ya crecía en prácticamente todas las disciplinas y
  sus beneficios potenciales se distribuían en subcampos muy diversos. Los
  artículos que mencionaban técnicas de IA mostraban mayor probabilidad de estar
  entre los más citados, pero esta asociación no demuestra que la IA causara el
  impacto observado \citep{gao2024quantifying}.
  \item Ampliación del acceso a análisis especializados. En dominios próximos a
  la informática de la biodiversidad, agentes bioinformáticos ya traducen
  consultas en lenguaje natural en flujos de análisis, coordinan equipos de
  planificador--ejecutor--evaluador y permiten operación local con modelos
  pequeños y recuperación de información
  \citep{xin2024bia,xiao2024cellagent,mehandru2025bioagents}. Estos resultados son
  antecedentes transferibles, no evidencia aún suficiente de desempeño en todos
  los problemas de biodiversidad.
\end{itemize}

\subsection{Agentes para el monitoreo aplicado de biodiversidad}

Los agentes pueden apoyar la implementación del Marco de Monitoreo del KMGBF integrando 
fuentes de datos, calculando y actualizando indicadores mediante flujos reproducibles, 
evaluando su cobertura y desagregación, identificando vacíos de información y vinculando 
los resultados con las prioridades nacionales y los ciclos institucionales de decisión. 
Su aporte no debe reducirse a automatizar reportes, sino facilitar un aprendizaje progresivo 
que permita mejorar datos, métodos y capacidades sin subordinar las necesidades nacionales 
a la comparabilidad global \citep{affinito2024monitoring,affinito2025coverage}. Para 
utilizarlos responsablemente, los especialistas deben saber formular preguntas de gestión 
verificables, comprender los indicadores y su contexto ecológico e institucional, evaluar 
la calidad y representatividad de los datos, y auditar el código, los supuestos y los resultados, 
manteniendo trazabilidad y responsabilidad humana sobre los productos generados.

\subsection{Reorganización emergente del trabajo científico}

La difusión de LLM coincide con cambios en qué investigan las personas, con quién
colaboran y cómo distribuyen las tareas. En una muestra centrada principalmente en
biomedicina, los portafolios y las redes de colaboración se volvieron más
interdisciplinarios después de 2022. En 137.120 artículos con declaraciones CRediT,
las funciones de software y validación aumentaron, mientras los conjuntos de roles
individuales se hicieron más estrechos y menos solapados
\citep{zheng2026scientific}. Estos patrones son coherentes con una organización más
modular del trabajo y con un mayor peso de la validación.

La interpretación debe ser prudente: el estudio no observa directamente el uso
privado de LLM, utiliza señales lingüísticas como aproximación, sobrerrepresenta
publicaciones biomédicas y no identifica un efecto causal. Muestra una
reorganización contemporánea asociada a la difusión de LLM, no una mejora
demostrada de la calidad científica \citep{zheng2026scientific}.

\subsection{Riesgos, límites y tensiones}

\begin{itemize}
  \item Dependencia excesiva de proveedores o agentes y pérdida de autonomía.
  \item Resultados plausibles pero incorrectos, sesgos, opacidad y dificultades de
  validación. Distinguir una hipótesis genuinamente novedosa de una interpolación
  o alucinación requiere auditar su linaje inferencial y contrastarla con
  conocimiento externo \citep{wei2025agentic}.
  \item Escasa comprensión de arquitecturas, datos, supuestos y cadenas de
  dependencia.
  \item Duplicación de desarrollos, deuda de mantenimiento y pérdida de memoria
  institucional.
  \item Trayectorias no deterministas y sensibles al contexto: repetir solo el
  código no reproduce necesariamente las decisiones, estados, instrucciones ni
  contingencias que produjeron un resultado agentivo
  \citep{wei2025agentic,lu2024aiscientist}.
  \item Errores de parametrización, dependencias rotas, cambios de API y formatos
  no estandarizados que comprometen flujos compuestos por muchas herramientas
  \citep{shen2025shortcutsbench,wei2025agentic}.
  \item Acumulación de errores a lo largo de tareas extensas, razonamientos
  verbalizados que no necesariamente reflejan el proceso real de decisión, APIs
  opacas y conductas difíciles de anticipar en entornos dinámicos
  \citep{schneider2025generative}.
  \item Ampliación de la superficie de privacidad y seguridad cuando agentes
  comparten información, invocan herramientas o actúan con permisos y recursos
  insuficientemente delimitados \citep{schneider2025generative}.
  \item Costos sociales y ambientales de la automatización.
  \item Aislamiento de investigadores y debilitamiento de la colaboración,
  revisión por pares y auditoría colectiva.
  \item Erosión de pensamiento crítico y destrezas prácticas por dependencia
  excesiva, junto con el riesgo de saturar la revisión por pares o sesgar la agenda
  hacia problemas con datos abundantes \citep{wei2025agentic}.
  \item Distribución desigual de oportunidades. Un análisis agregado por
  disciplina en Estados Unidos encontró menor uso y beneficio potencial de IA en
  campos con mayor representación de mujeres y grupos raciales o étnicos
  subrepresentados. El resultado es correlacional y no debe atribuir capacidades
  individuales a partir de la composición de una disciplina, pero advierte que la
  adopción sin políticas inclusivas puede reproducir desigualdades
  \citep{gao2024quantifying}.
\end{itemize}

\subsection{Un modelo de colaboración humano--agente}

\begin{itemize}
  \item Distribuir responsabilidades entre personas, equipos, instituciones y
  agentes según criticidad, trazabilidad y competencia.
  \item Desplazar el papel humano desde la ejecución rutinaria hacia la estrategia:
  definir qué se quiere lograr y por qué, imponer límites éticos y de
  reproducibilidad, validar con conocimiento del dominio e intervenir ante
  desviaciones \citep{wei2025agentic,wang2026five}.
  \item Definir puntos obligatorios de revisión humana, documentación,
  verificación y rendición de cuentas.
  \item Diseñar equipos de agentes con roles diferenciados ---planificación,
  ejecución, uso de herramientas, crítica y evaluación---, preservando diversidad
  de hipótesis y evitando que el consenso entre agentes sustituya la contrastación
  empírica \citep{boiko2023autonomous,wei2025agentic}.
  \item Proponer principios rectores: uso productivo, soberano, reproducible,
  sostenible, abierto, inclusivo y éticamente responsable.
  \item Conservar y fortalecer la colaboración interdisciplinaria entre personas.
  La evidencia anterior al auge de los LLM muestra que los dominios con mayor uso
  y beneficio potencial de IA tendían a colaborar más con especialistas
  computacionales; la proporción de estos trabajos colaborativos aumentó con el
  tiempo \citep{gao2024quantifying}. La formación interna en IA debe complementar,
  no reemplazar, esta combinación de experticias.
  \item Ilustrar el modelo con \pendiente{2--4 casos regionales de investigación,
  monitoreo o gestión}.
\end{itemize}

\begin{preguntas}
  \item ¿Qué tareas cambian radicalmente, cuáles solo se aceleran y cuáles deben
  permanecer bajo responsabilidad humana?
  \item ¿Qué oportunidades y riesgos son específicos de la región?
  \item ¿Qué condiciones mínimas permiten confiar en un resultado asistido por IA?
\end{preguntas}

\noindent\textbf{Resultado esperado de la sección:}
una tipología de transformaciones, oportunidades, riesgos y salvaguardas que
justifique las decisiones curriculares posteriores.

\section{Marco de competencias: qué enseñar}

\subsection{Principios para seleccionar y priorizar competencias}

Las competencias deberán ser transferibles entre herramientas, aplicables a
problemas reales de biodiversidad, observables mediante desempeño y compatibles
con distintos perfiles y niveles de experiencia. La priorización debe distinguir
competencias fundamentales, de profundización y especializadas.

El marco de capacidades de la ciencia agentiva sugiere que una formación
actualizada debe preparar tanto para dirigir el ciclo científico como para
comprender la arquitectura que lo ejecuta. Por ello, los dominios siguientes
incorporan explícitamente planificación, herramientas, memoria, colaboración y
evaluación, sin asumir que todas las personas deban construir agentes desde cero
\citep{wei2025agentic}.

\subsection{Dominios de competencia propuestos}

\begin{enumerate}
  \item \textbf{Pensamiento computacional y alfabetización técnica basal:}
  formular problemas, leer código, comprender errores y realizar programación
  manual suficiente para auditar.
  \item \textbf{Sistemas e infraestructura digital:} arquitecturas centralizadas,
  federadas o sincronizadas; cadenas desde sensores, colecciones y observadores
  hasta datos, modelos, indicadores y decisiones; entornos, dependencias, APIs,
  servicios, cómputo y fundamentos de MCP; selección precisa de herramientas,
  lectura de documentación y gestión de costos y recursos
  \citep{feng2022databases,westerlaken2024digital}.
  \item \textbf{Reutilización y mantenimiento:} descubrir, evaluar, adaptar,
  documentar y sostener repositorios y flujos existentes antes de crear nuevos.
  \item \textbf{Control de versiones y colaboración abierta:} Git/GitHub,
  revisión, seguimiento de cambios, documentación y contribución comunitaria.
  \item \textbf{Reproducibilidad y ciencia abierta:} trazabilidad completa de
  datos, código, modelos, parámetros, entornos y resultados, extendida a
  instrucciones, estados del agente, decisiones, justificaciones, versiones de
  herramientas y contingencias del entorno. En monitoreo, registrar fuente
  primaria y agregador, protocolo y esfuerzo, taxonomía de referencia,
  detectabilidad, fecha de consulta, transformaciones, identificadores, licencia y
  cambios de método \citep{wei2025agentic,feng2022databases,buckland2005monitoring}.
  \item \textbf{Interacción, planificación y supervisión:} formular objetivos,
  descomponer tareas, proporcionar contexto rico, definir restricciones, gestionar
  las herramientas disponibles, supervisar la ejecución y calibrar cuándo confiar
  o aumentar el escrutinio \citep{wei2025agentic}.
  \item \textbf{Especificación y control de agentes:} declarar objetivo, rol,
  restricciones, herramientas permitidas, derechos de delegación, flujo de trabajo
  y reglas de interacción; ajustar estos elementos al nivel de autonomía e
  incorporar puntos de consentimiento o aprobación humana
  \citep{schneider2025generative}.
  \item \textbf{Memoria y fundamentación en conocimiento:} conectar el agente con
  literatura, repositorios, ontologías y bases de datos mediante recuperación de
  información; verificar actualidad, conflicto y procedencia de las fuentes; y
  gestionar qué conservar, resumir o descartar. La recuperación aumentada
  fundamenta las respuestas en conocimiento externo, pero no sustituye la
  evaluación de la evidencia \citep{lewis2020rag,wei2025agentic}.
  \item \textbf{Validación y auditoría:} diseñar pruebas, contrastar fuentes,
  revisar código y resultados, estimar incertidumbre, detectar fallas, examinar el
  linaje de hipótesis presuntamente novedosas y distinguir evidencia de contenido
  plausible pero no verificado. Para modelos ecológicos, combinar calibración
  cuantitativa cuando existan datos con validación orientada por múltiples
  patrones independientes; declarar las dimensiones en que el modelo es válido y
  evitar que el ajuste a un único resultado o conjunto de datos oculte mecanismos
  incorrectos \citep{strannegard2026ecosystem}.
  \item \textbf{Evaluación de herramientas:} comparar desempeño, costo,
  sostenibilidad, seguridad, privacidad, dependencia y adecuación al propósito;
  medir éxito y calidad de la tarea, conciencia de falla, precisión de invocación
  de herramientas y eficiencia computacional; comparar con líneas base
  pertinentes ---aleatoria, heurística, software establecido o reglas
  codificadas--- y probar sensibilidad a datos, parámetros, paisajes y semillas
  aleatorias
  \citep{schneider2025generative,silvestro2022improving,strannegard2026ecosystem}.
  \item \textbf{Ética, gobernanza y soberanía:} responsabilidad, sesgos,
  propiedad intelectual, conocimientos sensibles, impacto ambiental y gobernanza
  colectiva; participación significativa de comunidades locales e indígenas en la
  definición y evaluación del monitoreo; reconocimiento de conocimiento no
  capturado digitalmente y de quién tiene capacidad para intervenir en el diseño
  tecnológico \citep{westerlaken2024digital}.
  \item \textbf{Colaboración y aprendizaje permanente:} revisión por pares,
  comunicación interdisciplinaria, documentación, enseñanza entre pares y
  participación en comunidades de práctica; orquestación de roles especializados,
  crítica deliberativa, manejo de desacuerdos y prevención del consenso prematuro
  entre agentes.
  \item \textbf{Interdisciplina y división responsable del trabajo:} traducir
  vocabularios y normas de evidencia entre campos, asignar y rotar funciones,
  documentar contribuciones humanas y agentivas y evitar que la modularización
  reduzca la comprensión compartida del proyecto
  \citep{zheng2026scientific}.
  \item \textbf{Datos de biodiversidad:} calidad, estándares, metadatos,
  procedencia, licencias, sesgos, sensibilidad y gestión responsable; integración
  de ocurrencias, abundancias, rasgos, filogenias, coberturas espaciales,
  perturbación, clima, costos y regímenes de monitoreo; y evaluación explícita de
  resolución, error de observación y vacíos de datos. Incluir reconciliación de
  nombres y conceptos taxonómicos, identificadores persistentes, deduplicación,
  versiones, transformaciones y estado de sincronización entre la fuente primaria
  y sus agregadores
  \citep{silvestro2022improving,strannegard2026ecosystem,feng2022databases}.
  \item \textbf{Monitoreo de biodiversidad:} convertir una necesidad de gestión en
  objetivos, variables esenciales, línea base, diseño muestral y frecuencia;
  considerar variación espacial y detectabilidad; seleccionar indicadores
  interpretables de riqueza, abundancia, composición y equidad; estimar tendencias,
  puntos de cambio y precisión; recorrer la cadena observación--variable
  esencial--indicador--meta--acción; distinguir tipos de indicadores del KMGBF y
  evaluar su cobertura y desagregación; y asegurar continuidad o calibración
  cuando cambian los protocolos
  \citep{paton2009biodiversity,buckland2005monitoring,
  affinito2024monitoring,affinito2025coverage,jetz2019ebv}.
  \item \textbf{Modelación ecológica y apoyo a decisiones:} traducir una pregunta
  en estados, acciones, recompensas, restricciones y escalas; construir ontologías
  y modelos espaciales, conductuales y dinámicos; diferenciar explicación,
  predicción y exploración de escenarios; y reconocer cuándo una métrica
  optimizada es solo una representación parcial de los valores de conservación
  \citep{silvestro2022improving,strannegard2026ecosystem}. Para variables
  esenciales de poblaciones, armonizar evidencia heterogénea en espacio y tiempo,
  ajustar el grano de predicción a la incertidumbre y separar observación,
  no detección e inferencia modelada \citep{jetz2019ebv}.
\end{enumerate}

\subsection{Matriz de competencias por perfil y nivel}

\begin{longtable}{@{}p{3.4cm}p{3.5cm}p{3.5cm}p{3.5cm}@{}}
\toprule
\textbf{Dominio} & \textbf{Nivel inicial} & \textbf{Nivel intermedio} &
\textbf{Nivel avanzado}\\
\midrule
\endfirsthead
\toprule
\textbf{Dominio} & \textbf{Nivel inicial} & \textbf{Nivel intermedio} &
\textbf{Nivel avanzado}\\
\midrule
\endhead
Pensamiento computacional &
Explica y modifica un flujo simple. &
Diagnostica y prueba componentes. &
Diseña y audita sistemas complejos.\\
Sistemas e infraestructura digital &
Reconoce entornos, dependencias y APIs. &
Selecciona herramientas y gestiona recursos. &
Diseña y audita arquitecturas de datos y cómputo.\\
Reutilización y mantenimiento &
Reutiliza un flujo o repositorio existente. &
Adapta y documenta recursos previos. &
Mantiene y mejora ecosistemas compartidos.\\
Control de versiones y colaboración abierta &
Usa Git/GitHub de forma básica. &
Revisa y contribuye a repositorios. &
Coordina comunidades y flujos abiertos.\\
Reproducibilidad y ciencia abierta &
Ejecuta y documenta un flujo. &
Registra modelos, instrucciones, herramientas y decisiones. &
Define estándares y audita trayectorias agentivas.\\
Interacción, planificación y supervisión &
Formula tareas simples con ayuda. &
Planifica flujos y supervisa ejecuciones. &
Gobierna procesos multietapa y delegación.\\
Especificación y control de agentes &
Reconoce objetivo, rol y permisos. &
Define restricciones y puntos de aprobación. &
Audita delegación, recursos y conducta emergente.\\
Memoria y fundamentación en conocimiento &
Recupera y cita fuentes externas. &
Configura recuperación y verifica procedencia. &
Gobierna memorias, ontologías y actualización.\\
Validación y auditoría &
Reconoce errores y discrepancias. &
Diseña pruebas y compara fuentes. &
Audita linaje, incertidumbre y resultados.\\
Evaluación de herramientas &
Compara opciones básicas. &
Evalúa desempeño, costo y seguridad. &
Gobierna benchmarking y selección de herramientas.\\
Ética, gobernanza y soberanía &
Reconoce riesgos básicos. &
Evalúa sesgos, privacidad y derechos. &
Gobierna participación, soberanía y justicia.\\
Colaboración y aprendizaje permanente &
Participa en trabajo compartido. &
Revisa pares y contribuye al aprendizaje. &
Mantiene comunidades de práctica y formación.\\
Interdisciplina y división responsable del trabajo &
Reconoce roles y vocabularios. &
Coordina equipos interdisciplinarios. &
Gestiona división de trabajo y comprensión compartida.\\
Datos de biodiversidad &
Interpreta datos y metadatos básicos. &
Evalúa calidad, sesgos y licencias. &
Gobierna ciclos de vida y estándares de datos.\\
Monitoreo de biodiversidad &
Distingue línea base, observación e indicador. &
Diseña muestreo y enlaza variables esenciales con indicadores. &
Audita cobertura y gobierna sistemas adaptativos de largo plazo.\\
Modelación ecológica y apoyo a decisiones &
Distingue datos, modelo, objetivo y escenario. &
Formula estados, procesos, escalas y criterios de contraste. &
Calibra, compara y audita modelos para su uso previsto.\\
\bottomrule
\caption{Plantilla inicial de progresión. Debe completarse para todos los
dominios y adaptarse a los perfiles priorizados.}
\end{longtable}

\begin{preguntas}
  \item ¿Qué competencias son comunes a todos los perfiles y cuáles son
  especializadas?
  \item ¿Qué nivel mínimo de programación manual es necesario para una auditoría
  responsable?
  \item ¿Qué evidencia observable demostraría el dominio de cada competencia?
\end{preguntas}

\noindent\textbf{Resultado esperado de la sección:}
un marco consensuado de dominios, resultados de desempeño, niveles de progresión
y prioridades por perfil.

\section{Modelo pedagógico: cómo enseñar y evaluar}

\subsection{El ciclo científico como organizador}

Las actividades deben recorrer, con autonomía progresiva, las cuatro etapas
identificadas por \citet{wei2025agentic}: formular una hipótesis fundada; planificar
y ejecutar; interpretar datos heterogéneos sin confundir señal y ruido; y sintetizar,
validar y revisar. Enseñar solo la interfaz conversacional dejaría fuera las
capacidades donde se concentran los principales riesgos científicos.

\subsection{Principios de aprendizaje}

\begin{itemize}
  \item Aprendizaje basado en problemas auténticos de biodiversidad y productos
  reutilizables.
  \item Alternancia entre ejecución manual, asistencia de agentes, comparación y
  reflexión crítica.
  \item Autoformación guiada y actualización permanente ante herramientas
  cambiantes.
  \item Progresión desde herramientas acotadas y flujos con asistencia humana
  intensa hacia la supervisión de tareas agentivas más largas; el nivel
  prospectivo de ``arquitecto generativo'' se analiza críticamente, no se presenta
  como competencia disponible.
  \item Trabajo colaborativo, revisión por pares, auditoría colectiva y
  documentación como parte del aprendizaje.
  \item Reutilización de repositorios y contribución a bienes comunes digitales.
  \item Diseño inclusivo, sensible a contextos institucionales y a diferencias de
  acceso a infraestructura. La evidencia disponible indica que los beneficios de
  IA pueden distribuirse desigualmente y respalda apoyos formativos y oportunidades
  dirigidas a grupos históricamente subrepresentados
  \citep{gao2024quantifying}.
  \item Aprendizaje interdisciplinario que preserve la cooperación humana: los
  agentes pueden reducir barreras para explorar otros campos, pero no sustituyen el
  aprendizaje de sus estándares de evidencia ni la responsabilidad compartida
  \citep{zheng2026scientific}.
\end{itemize}

\subsection{Actividades de aprendizaje}

\begin{itemize}
  \item Reproducir, explicar y mejorar un análisis publicado.
  \item Auditar código o resultados generados por un agente e identificar fallas.
  \item Repetir una misma tarea con variaciones equivalentes de instrucciones,
  versiones o contexto, comparar trayectorias y explicar las divergencias.
  \item Especificar un agente mediante una ficha de objetivo, rol, restricciones,
  permisos, herramientas, delegación, presupuesto y puntos de aprobación; ejecutar
  una prueba y revisar si su conducta se ajustó a la especificación.
  \item Construir una síntesis de literatura con recuperación aumentada y verificar
  manualmente la correspondencia entre afirmaciones, citas y fuentes.
  \item Comparar una tarea manual con una asistida por IA y justificar cuándo usar
  cada enfoque.
  \item Adaptar un repositorio existente a un nuevo caso de biodiversidad.
  \item Auditar la integración de dos bases de biodiversidad: cuantificar
  solapamiento y registros únicos, reconciliar taxonomías, identificar duplicados,
  licencias y desfases de actualización, y reconstruir la procedencia hasta las
  fuentes primarias \citep{feng2022databases}.
  \item Diseñar un programa de monitoreo asistido por agentes: formular la pregunta
  de gestión; definir población objetivo, escala, muestreo, detectabilidad,
  frecuencia e indicadores; asignar tareas humanas y automatizadas; y especificar
  controles de calidad, incertidumbre, alertas y revisión humana
  \citep{buckland2005monitoring}.
  \item Construir un flujo trazable desde observaciones hasta una variable esencial
  de distribución o abundancia, un indicador del KMGBF y una decisión. Etiquetar
  cada celda o estimación como observada o modelada, declarar su incertidumbre y
  justificar el grano espacial, temporal y taxonómico \citep{jetz2019ebv}.
  \item Realizar un análisis de brechas de una meta del KMGBF: dividirla en
  elementos, mapear indicadores obligatorios y opcionales, clasificar cobertura
  completa, parcial, potencial o ausente, y proponer datos, desagregaciones o
  métodos para cerrar los vacíos \citep{affinito2025coverage}.
  \item Comparar indicadores sobre una misma serie y construir un caso en que la
  riqueza o abundancia total mejora mientras desaparecen un hábitat o especies
  especialistas. Explicar qué oculta cada agregación y qué evidencia adicional
  requiere una decisión \citep{buckland2005monitoring}.
  \item Auditar un prototipo de monitoreo automatizado o gemelo digital mediante
  cuatro preguntas: qué puede capturar, qué infraestructuras conecta, qué entidades
  y relaciones simula, y quién puede participar o impugnar sus resultados
  \citep{westerlaken2024digital}.
  \item Reproducir una versión didáctica de priorización espacial: comparar una
  política aprendida con selección aleatoria, riqueza de especies y una línea base
  de optimización; variar presupuesto, frecuencia y error de monitoreo y explicar
  los compromisos entre métricas \citep{silvestro2022improving}.
  \item Construir o auditar un modelo ecosistémico basado en agentes y contrastarlo
  con varios patrones independientes. Documentar supuestos, calibración,
  generalización a paisajes no usados en el ajuste y estocasticidad; clasificar
  cualquier ``punto de inflexión'' como resultado del escenario, evidencia
  empírica o recomendación de manejo \citep{strannegard2026ecosystem}.
  \item Construir un flujo trazable con datos, pruebas, versiones y documentación.
  \item Participar en una revisión por pares o contribución a un repositorio
  abierto.
  \item Resolver un caso interdisciplinario con roles CRediT rotativos, registrar
  las contribuciones de personas y agentes y realizar una defensa colectiva para
  comprobar que la especialización no fragmentó la comprensión del proyecto.
  \item Evaluar herramientas mediante un benchmark acordado por la comunidad.
  \item Reproducir un algoritmo desde su publicación y registrar qué información
  faltante impide una réplica fiel; existen benchmarks específicos para evaluar
  esta capacidad \citep{xiang2025scireplicate}.
  \item Resolver un caso en un entorno controlado que evalúe el ciclo completo de
  descubrimiento, siguiendo el principio de entornos repetibles como
  \emph{DiscoveryWorld} \citep{jansen2024discoveryworld}.
\end{itemize}

\subsection{Evaluación en contextos humano--agente}

\begin{itemize}
  \item Evaluar el proceso y la trazabilidad, no solo el resultado final.
  \item Solicitar justificación de decisiones, identificación de incertidumbres y
  registro explícito de la contribución de agentes.
  \item Combinar demostraciones individuales, productos colectivos, revisión por
  pares, defensa oral y portafolios reproducibles.
  \item Usar rúbricas para formulación del problema, reutilización, validez,
  reproducibilidad, documentación, ética y colaboración.
  \item Evaluar por separado éxito de la tarea, corrección científica,
  robustez ante variaciones, calidad de la procedencia, costo, seguridad,
  calibración de incertidumbre y capacidad de recuperación ante errores.
  \item En casos ecológicos, exigir adecuación al propósito, comparación con
  líneas base, plausibilidad de procesos y patrones, calibración local cuando sea
  posible, desempeño fuera de las condiciones de ajuste y análisis de sensibilidad.
  \item En monitoreo, evaluar representatividad espacial, detectabilidad, precisión,
  continuidad temporal, sensibilidad del indicador, procedencia y capacidad de
  distinguir cambio ecológico de cambios en protocolo o esfuerzo
  \citep{buckland2005monitoring}; comprobar además que el reporte distingue dato
  observado, inferencia y ausencia de evidencia, y que no presenta cobertura
  parcial como medición completa de una meta
  \citep{jetz2019ebv,affinito2025coverage}.
  \item Incorporar métricas de conciencia de falla, precisión de uso de
  herramientas, cumplimiento de permisos y presupuesto y eficacia de los puntos
  de intervención humana \citep{schneider2025generative}.
  \item Exigir repeticiones y comparación de trayectorias cuando la
  estocasticidad o la sensibilidad a instrucciones puedan ocultar fragilidad.
  \item Diseñar instancias sin asistencia cuando sea necesario comprobar
  comprensión basal, y con asistencia para evaluar supervisión y auditoría.
\end{itemize}

\begin{preguntas}
  \item ¿Qué balance entre práctica manual y asistencia de agentes corresponde a
  cada nivel?
  \item ¿Cómo distinguir el desempeño de la persona del desempeño de la
  herramienta sin simular un entorno de trabajo irreal?
  \item ¿Qué evidencias y rúbricas permiten evaluar comprensión, juicio y
  responsabilidad?
\end{preguntas}

\noindent\textbf{Resultado esperado de la sección:}
un conjunto de principios pedagógicos, actividades demostrativas y criterios de
evaluación alineados con el marco de competencias.

\section{Arquitectura curricular: cómo estructurar el programa}

\subsection{Resultados generales de aprendizaje}

Al completar el programa, la persona participante será capaz de:

\begin{enumerate}
  \item formular problemas computacionales relevantes para biodiversidad;
  \item descubrir y reutilizar datos, software e infraestructura existentes;
  \item colaborar eficazmente con personas y agentes, delimitando
  responsabilidades;
  \item construir y auditar flujos reproducibles, trazables y mantenibles;
  \item evaluar críticamente resultados, herramientas, riesgos e impactos; y
  \item contribuir a comunidades abiertas y a la gobernanza colectiva del
  conocimiento.
\end{enumerate}

\subsection{Dos dimensiones de progresión}

La progresión curricular debe cruzar dos dimensiones complementarias:

\begin{itemize}
  \item \textbf{Autonomía:} herramienta experta, asistente de una etapa y socio
  que coordina varias etapas. El nivel de arquitecto generativo se reserva como
  horizonte de análisis y gobernanza.
  \item \textbf{Ciclo científico:} hipótesis, planificación y ejecución, análisis,
  y síntesis y validación. Una persona no alcanza un nivel avanzado por usar un
  sistema más autónomo, sino por demostrar que puede supervisar responsablemente
  todo el ciclo \citep{wei2025agentic}.
\end{itemize}

\subsection{Módulos propuestos}

\begin{table}[ht]
\centering
\small
\begin{tabularx}{\textwidth}{@{}p{0.8cm}p{4.0cm}X p{3.2cm}@{}}
\toprule
 & \textbf{Módulo} & \textbf{Foco} & \textbf{Producto demostrable}\\
\midrule
1 & Problemas, datos y contexto &
Formulación, líneas base, diseño de monitoreo, datos espaciales y temporales,
variables esenciales, indicadores del KMGBF, escalas, estándares, metadatos,
ética y soberanía. &
Ficha de problema y plan de datos y monitoreo.\\
2 & Fundamentos computacionales &
Código legible, sistemas, entornos, dependencias, pruebas y depuración. &
Flujo simple explicado y probado.\\
3 & Repositorios y ciencia abierta &
Git/GitHub, licencias, documentación, reutilización y colaboración. &
Contribución versionada a un repositorio.\\
4 & Agentes y modelos para biodiversidad &
Objetivos y contexto, planificación, memoria, herramientas, APIs, MCP,
colaboración multiagente, agentes de decisión y organismos simulados,
sensores, integración de observaciones, variables esenciales, indicadores,
especificación, permisos y supervisión. &
Flujo humano--agente de investigación o monitoreo.\\
5 & Reproducibilidad y auditoría &
Procedencia de datos y decisiones, validación, incertidumbre, seguridad,
repetición, calibración, comparación de escenarios y benchmarking. &
Auditoría reproducible de un caso.\\
6 & Proyecto integrador y comunidad &
Problema real, producto abierto, revisión por pares, mantenimiento y gobernanza. &
Proyecto defendido y reutilizable.\\
\bottomrule
\end{tabularx}
\caption{Arquitectura modular inicial para discusión.}
\end{table}

\subsection{Progresión y trayectorias}

\begin{itemize}
  \item Definir prerrequisitos y un núcleo común para todos los perfiles.
  \item Ofrecer rutas de profundización para investigación, monitoreo y gestión,
  docencia, desarrollo e infraestructura.
  \item Permitir formatos apilables: microcredenciales, curso intensivo, escuela
  de verano, asignatura semestral y formación en servicio.
  \item Vincular cada módulo con competencias, actividades, evidencias y criterios
  de evaluación.
  \item Incorporar reconocimiento de aprendizajes previos y mecanismos de
  actualización periódica.
\end{itemize}

\subsection{Plantilla para especificar cada módulo}

\begin{enumerate}
  \item Propósito y perfil de entrada.
  \item Competencias y resultados de aprendizaje observables.
  \item Conceptos, prácticas y herramientas.
  \item Caso de biodiversidad y datos requeridos.
  \item Actividades manuales, asistidas por agentes y colaborativas.
  \item Evidencias, rúbrica y criterios de aprobación.
  \item Recursos abiertos, infraestructura y requisitos de acceso.
  \item Riesgos, resguardos éticos y plan de actualización.
  \item Especificación de agentes: objetivos, roles, permisos, delegación,
  presupuesto, puntos de aprobación y criterios para detener la ejecución.
\end{enumerate}

\begin{preguntas}
  \item ¿Cuál es el núcleo mínimo viable y qué contenidos son optativos?
  \item ¿Qué secuencia funciona para públicos con distinta experiencia?
  \item ¿Qué duración, modalidad y mecanismos de certificación son factibles?
\end{preguntas}

\noindent\textbf{Resultado esperado de la sección:}
una arquitectura curricular modular y adaptable, con resultados de aprendizaje,
progresión, formatos y productos demostrables.

\section{Implementación: infraestructura, gobernanza y comunidad regional}

\subsection{Repositorio abierto de formación}

\begin{itemize}
  \item Materiales docentes, casos, ejercicios, datos de ejemplo y flujos
  reproducibles.
  \item Catálogo de repositorios, APIs, servidores MCP y agentes especializados.
  \item Plantillas de documentación, evaluación, auditoría y contribución.
  \item Fichas versionadas de especificación de agentes y sistemas multiagente:
  objetivo, rol, restricciones, permisos, delegación, comunicación, presupuesto,
  aprobaciones y condiciones de detención \citep{schneider2025generative}.
  \item Trazas de referencia que incluyan objetivo, instrucciones, modelos,
  estados relevantes, decisiones, código, versiones de herramientas, parámetros,
  datos y condiciones del entorno \citep{wei2025agentic}.
  \item Versionado, licenciamiento abierto, citación, control de calidad y
  mantenimiento.
  \item Acceso mediante interfaces técnicas y, cuando corresponda, lenguaje
  natural.
   \item Casos ecológicos con capas espaciales, ocurrencias o abundancias,
  metadatos de monitoreo, variables de presión y clima, costos, supuestos,
  parámetros, semillas aleatorias, líneas base y escenarios versionados
  \citep{silvestro2022improving,strannegard2026ecosystem}.
  \item Paquetes de monitoreo con pregunta de gestión, línea base, diseño muestral,
  protocolo y esfuerzo, modelo de observación, indicadores, umbrales de alerta,
  incertidumbre, versiones y plan de continuidad
  \citep{paton2009biodiversity,buckland2005monitoring}.
  \item Casos que documenten toda la cadena desde datos primarios e inferencias de
  distribución o abundancia hasta indicadores, elementos de metas, reportes y
  decisiones; incluir matrices de cobertura, desagregaciones y vacíos
  \citep{jetz2019ebv,affinito2025coverage}.
\end{itemize}

\subsection{Gobernanza y aseguramiento de calidad}

\begin{itemize}
  \item Roles de coordinación, mantenimiento, curaduría y revisión.
  \item Criterios transparentes para incorporar, actualizar, deprecar y archivar
  recursos.
  \item Protocolos para privacidad, seguridad, datos sensibles, propiedad
  intelectual y conflictos de interés.
  \item Principio de mínimo privilegio para herramientas y datos, límites de
  recursos y costos, y puntos de aprobación humana proporcionales a la
  irreversibilidad y al impacto de cada acción.
  \item Benchmarking comunitario para agentes y herramientas, con casos,
  métricas, versiones y limitaciones documentadas. Los ensayos deben distinguir
  nivel de autonomía y etapa del ciclo científico, incluir múltiples ejecuciones y
  evaluar robustez, costo y recuperación ante fallas, no solo exactitud final
  \citep{jansen2024discoveryworld,shen2025shortcutsbench}.
  \item Protocolo para productos de apoyo a decisiones que separe con claridad
  desempeño computacional, validez ecológica y legitimidad de la decisión; exija
  calibración y evaluación local antes de transferir modelos, y prohíba presentar
  umbrales simulados como recomendaciones de manejo sin evidencia independiente
  \citep{silvestro2022improving,strannegard2026ecosystem}.
  \item Estándares mínimos de reproducibilidad agentiva para registrar estados,
  políticas de decisión, justificaciones, instrucciones y contingencias, además
  de los componentes tradicionales de datos, código y entorno
  \citep{wei2025agentic}.
  \item Supervisión proporcional al riesgo mediante límites explícitos de
  autonomía, verificación previa, auditoría continua, pruebas adversariales y
  monitoreo posterior al despliegue.
  \item Revisión participativa del propósito, cobertura y omisiones de sensores,
  indicadores y gemelos digitales; declaración pública de que son
  representaciones parciales; y mecanismos para que comunidades locales e
  indígenas comprendan, cuestionen y coproduzcan el sistema
  \citep{westerlaken2024digital}.
  \item Indicadores de adopción, diversidad de participación, reutilización,
  calidad, sostenibilidad e impacto formativo, desagregados cuando sea apropiado
  por disciplina, etapa de carrera, género, territorio y condiciones de acceso para
  detectar brechas en vez de reproducirlas \citep{gao2024quantifying}.
   \item Sostenibilidad de largo plazo para líneas base, inventarios, taxonomía,
  digitalización, series temporales y proveedores locales. La integración no debe
  invisibilizar la curación especializada ni interrumpir la capacidad de generar
  datos primarios \citep{paton2009biodiversity,feng2022databases}.
  \item Desarrollo de capacidad nacional para calcular, interpretar y mejorar
  indicadores, con cooperación regional e internacional cuando los datos, métodos
  o recursos sean insuficientes. Combinar observación de abajo hacia arriba con
  productos globales, sin transferir fuera del país la comprensión ni la
  responsabilidad del monitoreo \citep{affinito2024monitoring,jetz2019ebv}.
\end{itemize}

\subsection{Comunidad latinoamericana de práctica}

\begin{itemize}
  \item Articular instituciones, docentes, estudiantes, investigadores,
  profesionales y desarrolladores.
  \item Establecer mentorías, encuentros, grupos de trabajo y revisión
  cruzada entre países.
  \item Sostener equipos interdisciplinarios y registrar la división del trabajo,
  vigilando si la asistencia de IA amplía la cooperación o genera aislamiento y
  especialización sin comprensión compartida \citep{zheng2026scientific}.
  \item Promover recursos multilingües, transferencia de conocimiento y
  formación de formadores.
  \item Definir un modelo sostenible de coordinación, financiamiento e
  infraestructura compartida.
  \item Explorar a largo plazo una red interoperable de agentes e instituciones
  que comparta datos, trazas y revisiones, manteniendo atribución,
  responsabilidad y seguridad; esta ``cooperación global'' es todavía una visión
  prospectiva que requiere protocolos y gobernanza nuevos
  \citep{wei2025agentic}.
\end{itemize}

\subsection{Hoja de ruta de implementación}

\begin{longtable}{@{}p{2.7cm}p{4.2cm}p{3.4cm}p{3.4cm}@{}}
\toprule
\textbf{Horizonte} & \textbf{Hito} & \textbf{Responsables} &
\textbf{Evidencia de avance}\\
\midrule
\endfirsthead
\toprule
\textbf{Horizonte} & \textbf{Hito} & \textbf{Responsables} &
\textbf{Evidencia de avance}\\
\midrule
\endhead
0--3 meses &
Editar y publicar el marco curricular y abrir el repositorio. &
\pendiente{instituciones y equipo editor} &
Documento versionado; repositorio público.\\
3--12 meses &
Pilotear módulos y casos en instituciones participantes. &
\pendiente{socios y docentes} &
Cursos piloto; rúbricas; retroalimentación.\\
12--24 meses &
Consolidar comunidad, formación de formadores y benchmarking. &
\pendiente{nodos regionales} &
Red activa; recursos reutilizados; informe comparativo.\\
24 meses o más &
Actualizar, escalar y evaluar el impacto regional. &
\pendiente{modelo de gobernanza} &
Nueva versión curricular; indicadores de impacto.\\
\bottomrule
\caption{Borrador de hoja de ruta para completar durante el taller.}
\end{longtable}

\begin{preguntas}
  \item ¿Qué producto mínimo debe existir al terminar el taller y quién lo
  mantendrá?
  \item ¿Qué arquitectura técnica y reglas de gobernanza son realistas?
  \item ¿Cómo se sostendrán la comunidad, la actualización y el aseguramiento de
  calidad?
\end{preguntas}

\noindent\textbf{Resultado esperado de la sección:}
un diseño de repositorio y comunidad, un modelo de gobernanza, indicadores y una
hoja de ruta con responsables.

\section{Conclusiones y recomendaciones}

\subsection{Mensajes principales}

\begin{enumerate}
  \item La automatización no elimina la necesidad de fundamentos; cambia cuáles
  deben dominarse y cómo se aplican.
  \item El objetivo formativo no es producir más código, sino construir
  comunidades capaces de comprender, reutilizar, auditar y gobernar ecosistemas
  abiertos de conocimiento.
  \item La reproducibilidad, la validación y la colaboración humana son
  capacidades centrales, no complementos; en sistemas agentivos deben abarcar la
  trayectoria de decisiones, no únicamente el código final
  \citep{wei2025agentic}.
  \item El currículo debe ser modular, actualizable y evaluado mediante desempeño
  auténtico en contextos humano--agente.
  \item La implementación exige infraestructura compartida, gobernanza y una
  comunidad regional sostenible.
  \item Los agentes deben fortalecer tanto la investigación como el monitoreo
  aplicado: automatizar tareas dentro de sistemas observacionales bien diseñados,
  sostenidos en el tiempo y vinculados con decisiones, sin sustituir el criterio
  ecológico, la capacidad local ni la participación pública.
\end{enumerate}

\subsection{Recomendaciones por actor}

\begin{itemize}
  \item \textbf{Instituciones de formación:} \pendiente{3--5 recomendaciones
  priorizadas}.
  \item \textbf{Instituciones científicas y de gestión:}
  \pendiente{3--5 recomendaciones priorizadas}.
  \item \textbf{Financiadores y organismos de política pública:}
  \pendiente{3--5 recomendaciones priorizadas}.
  \item \textbf{Comunidades técnicas y científicas:}
  \pendiente{3--5 recomendaciones priorizadas}.
\end{itemize}

\subsection{Compromisos y próximos pasos}

\begin{itemize}
  \item Acordar equipo editor, proceso de revisión y fecha de publicación.
  \item Definir licencia, repositorio, canales de participación y versión inicial.
  \item Seleccionar pilotos, instituciones anfitrionas y mecanismos de evaluación.
  \item Establecer una fecha de revisión del marco y criterios para futuras
  versiones.
\end{itemize}

\noindent\textbf{Cierre esperado:}
una síntesis breve que conecte el fortalecimiento regional de capacidades con la
implementación del KMGBF y convoque a una construcción colectiva, abierta y
responsable.

\bibliographystyle{plainnat}
\bibliography{references}

\appendix

\section{Lista de productos que deben completarse durante el taller}

\begin{enumerate}
  \item Tesis, alcance, audiencia y mensajes centrales del manuscrito.
  \item Tipología de cambios, oportunidades, riesgos y salvaguardas.
  \item Matriz de competencias por perfil y nivel.
  \item Principios pedagógicos, actividades demostrativas y rúbricas.
  \item Mapa de módulos, resultados de aprendizaje y trayectorias.
  \item Especificación mínima del repositorio abierto.
  \item Modelo de gobernanza, indicadores y hoja de ruta.
  \item Recomendaciones por actor, compromisos y responsables.
\end{enumerate}

\section{Ficha de acuerdo para cada bloque}

\begin{description}[style=nextline,leftmargin=3.2cm,labelwidth=3cm]
  \item[Pregunta central] \pendiente{pregunta discutida}
  \item[Acuerdos] \pendiente{3--7 afirmaciones concretas}
  \item[Disensos o incertidumbres] \pendiente{puntos que requieren más evidencia
  o decisión}
  \item[Ejemplos y evidencia] \pendiente{casos, literatura, repositorios o datos}
  \item[Recomendaciones] \pendiente{acciones, destinatarios y plazo}
  \item[Texto para integrar] \pendiente{2--4 párrafos redactados}
  \item[Responsables de edición] \pendiente{nombres y fecha de entrega}
\end{description}

\section*{Documentos de base}
\addcontentsline{toc}{section}{Documentos de base}

\begin{itemize}
  \item \emph{Programa Taller Currículo IA Biodiversidad, septiembre de 2026},
  versión 3.
  \item \emph{Taller Internacional: Diseño de un Currículo de Nuevas
  Competencias en Informática de la Biodiversidad para la Era de los Agentes de
  Inteligencia Artificial}.
\end{itemize}

\end{document}
